\section{Fhara --- ``The Oasis Courts''}
\label{chap:fhara}

\subsection*{Quick Reference}
\textbf{Tagline:} \textit{Water is a gift from the ancestors. The well remembers every cup.}

\begin{quote}
  \textbf{Well-Judge (Elite):}  
  ``Water is not a right. It is a gift from the ancestors, passed down through the sand. The oasis remembers every drop drawn, every cup poured, every promise made at its edge. I do not judge the living. I witness the balance.''
\end{quote}

\begin{quote}
  \textbf{Water-Clerk (Commoner):}  
  ``My grandfather dug this well with a shovel and a prayer. My mother sealed the walls with clay from the dry river. I sit in the shade and measure every bucket. The sand does not forgive a broken trust. But it does forgive a forgetful mind – if you return to share.''
\end{quote}

\textbf{Genre / Mood:} Desert law, coffee ritual, water-as-community.

\textbf{Starting Location:} The shaded stone bench of a well-court, where the air is cool and the water is sweet. A Well-Judge dips a clay cup into the basin: ``You have come for water, or for trade, or for a story. First, drink. Then tell me which promise you intend to keep today.''

\vspace{2mm}
\begin{tcolorbox}[colback=black!3,colframe=blue!60!black,title={Theme \& Atmosphere}]
The Fhara dwell on a great, sun-baked peninsula where scattered oases and dry riverbeds connect the coast to the interior. Their cities are built around wells and date groves, and their caravans are the arteries of the spice trade. Water is law – the well-keepers are judges, and the courthouse is the shade of a palm. The Fhara are merchants, poets, and negotiators; they will bargain for an hour over a cup of coffee, and a day over a camel. Their honour is measured in generosity, and their memory is measured in water-shares. Coffee is not a drink; it is a ritual. Three cups: the first is welcome, the second is negotiation, the third is a contract. Refuse the third, and you have broken faith.

\vspace{2mm}
\textbf{What you notice first:} The smell of roasting coffee and overripe dates. The way every well has a stone bench for the judge. The silence when the sun passes directly overhead – no one moves, no one speaks.

\textbf{The one rule that kills newcomers:} Never refuse a date offered by a well-keeper. The fruit is sacred to their well-spirits. Refusing it ends the conversation – and your access to water.
\end{tcolorbox}

\subsection*{Regional Mood: The Weight of Water}

Power here is measured in well-depth, date groves, and the length of a coffee ritual. Each oasis is a sovereign court, ruled by a Sheikh and a Well-Judge. The Fhara have no single king; their confederacy is a web of marriage alliances, water-sharing treaties, and caravan tariffs. The desert is not empty – it is a record of dry wells, abandoned caravanserais, and the bones of those who could not share. A Fhara merchant can cross the waste with only a waterskin and a coffee pot, because every well has a keeper, every keeper a price, and every price can be negotiated – but the negotiation is a ritual of reciprocity, not of obligation.

\subsection*{The Antagonists: The Nerethi}
\label{subsec:nerethi}

From the deep desert comes the \textbf{Nerethi}. They are acentric, psionic, and distrusted by all. They do not believe in debt; to them, the concept is alien. They do not owe, they do not remember, and they do not forgive. They walk the dunes like shadows, their minds linked in a hive that rejects the Fhara's water-ledger.

\begin{itemize}
  \item \textbf{Philosophy:} Memory is a chain. Debt is a shackle. The Nerethi burn both.
  \item \textbf{Mechanic:} Nerethi agents can erase a \textbf{Bond} or \textbf{Obligation} from a character sheet, but the character loses the memory associated with it permanently.
  \item \textbf{Conflict:} The Fhara see them as thieves of history. The Nerethi see the Fhara as slaves to the past.
\end{itemize}

\subsection*{The Lost Capital: Qalat al-Madina (The Sunken City)}
\label{subsec:qalat}

Before the Kuvani and the Eastern Ykrul swept from the steppe, the Fhara had a single great city – Qalat al-Madina, the Mother of Wells. It stood where the peninsula's largest aquifer surfaced, a ring of gardens and fountains, its walls white as salt. Merchants from Kahfagia and Thepyrgos paid the city's Water-Tithe to cross the waste. Poets sang of its nine hundred wells, each with a different taste.

But Qalat al-Madina's true secret was older than the city itself. Beneath its deepest well, a faction of scholars unearthed a vault of \textbf{forbidden knowledge} – fragments of a civilisation that had turned to ash twenty thousand years before the first Utaran legion marched. The knowledge was not magic; it was a \textit{reckoning}. It spoke of a fire that could scour the world clean.

The three Omens – \textbf{Inaea, Ikasha, and Isoka} – appeared to the Eastern Ykrul and the Kuvani as whispers in the grass and cracks in the sky. They gave a single command: \textit{Bury it. Seal it. If it wakes, the world burns.} The horse-lords rode against Qalat al-Madina, not for plunder, but to honour the Omens' warning. They poisoned the wells, drowned the city in its own aquifer, and salted the earth. The curse they carved into the central well was not cruelty – it was a lock: \textit{``Let no Fhara ever drink here without paying a death.''} The death was the knowledge, buried alive.

Now Qalat al-Madina is a sinkhole field, half-submerged in brackish water, its towers leaning into the mud. The ghosts of the drowned well-keepers rise at dusk, still trying to measure water that no longer exists – and to guard the vault that no one has found. The Kuvani and Ykrul have long since left, but the \textbf{Omen-Whisper} remains: in certain wells, at certain tides, a traveller may hear a voice speaking a single word: \textit{``Remember.''}

\paragraph{What remains:}
\begin{itemize}
  \item \textbf{The Sunken Archive} – a flooded chamber where clay tablets still float, their judgments still legally binding on the living.
  \item \textbf{The Nine Hundred Wells} – now sinkholes, some connected to the underworld.
  \item \textbf{The Cursed Spring} – a pool that rises every new moon; those who drink from it forget their names for a year and a day.
  \item \textbf{The Drowned Coin} – a silver piece minted in the city's last year; it is said to buy safe passage from any Fhara well-spirit, because it still bears the Mother of Well's seal – and the memory of the Omens' warning.
\end{itemize}

\paragraph{Adventure Hooks:}
\begin{itemize}
  \item A Fhara sheikh wants to recover a water-tablet from the Sunken Archive. The tablet proves his clan never ceded the oasis – but the archive is haunted by the drowned judges who whisper of the forbidden vault.
  \item The Kuvani are sending an expedition to the ruins to \textit{renew} the curse. The Fhara need someone to sabotage them – or to learn why the Omens' seal is cracking.
  \item A ghost from Qalat al-Madina appears at a well-court, demanding justice for a murder committed three centuries ago. The accused is a modern Fhara noble, but the ghost's testimony is legally binding under ancient water law. The ghost also whispers a fragment of the forbidden knowledge: a single number.
  \item \textbf{Nerethi Incursion:} A Nerethi psion offers to erase the curse on Qalat al-Madina. The price: the party must forget why they came to the desert in the first place.
\end{itemize}

\begin{tcolorbox}[colback=black!2,colframe=red!70!black,title={The Last Scribe of Qalat al-Madina}]
  \textit{``I wrote what I should not have. Then I ran.''}
  
  The night the Kuvani and Eastern Ykrul rode against Qalat al-Madina, a young water-clerk named \textbf{Samir al-Qalat} sat in the deep archive, copying ancient tablets by lamplight. He was not a hero. He was the grandson of a well-keeper, and he had a bad habit of reading the tablets he was supposed to copy.
  
  That evening, he unrolled a scroll sealed with a symbol he did not recognise – three interlocking rings, each a different shade of green. The ink was still wet. The scroll had been written that morning, by a hand that no longer existed. It was the \textbf{Overseer's fragment}, smuggled out of a vault that should have stayed sealed.
  
  Samir read eight lines. Then he stopped. He knew that if he read the ninth, he would never stop.
  
  When the first torches appeared at the city's edge, he wrapped the scroll in oilcloth, stuffed it into a clay water-jar, and ran. Behind him, the wells were poisoned. Ahead, the desert. He walked for nine days, drinking only the moisture he could lick from the clay.
  
  On the ninth day, he stumbled into a Tulkani caravan circle. The Caravan Judge, an old woman with a cord that dragged on the ground, looked at the jar and said, \textit{``You have brought the ninth cup. We do not pour it. But we can carry it.''}
  
  Samir gave the scroll to the Tulkani. He did not ask for payment. He asked only that they keep it moving – that no single camp hold it for more than nine nights, no single wagon keep it sealed for more than a season. \textit{``The record is warm,''} he said. \textit{``The ink is still wet. The Hollow reads over my shoulder.''}
  
  The Tulkani honoured his request. The fragment passed from wagon to wagon, from story to story, until it reached the Archivolt of Thepyrgos, where it was sealed in a lead-lined box, encased in \textit{karadasarythm}, with a lock that requires eight witnesses to open. The ninth witness, as the \textit{Words of the Ninth} say, is always the Hollow – and the Hollow is always available.
  
  Samir al-Qalat died in the desert three years later, alone except for a single date-palm that grew from a seed he had carried from his mother's grove. The tree still stands. Travellers who rest in its shade sometimes hear a whisper: \textit{``Count to eight. Then stop.''}
  
  \vspace{2mm}
  \textbf{GM Hook:} The Tulkani still carry the memory of Samir's scroll. A Caravan Judge may offer the party a single glimpse of a copy – but only if they promise to carry a message to the Sunken City: \textit{``The ink is still wet. Has the Overseer finished writing?''}
  \end{tcolorbox}

\subsection*{Prominent NPCs of the Fhara}
\label{sec:fhara-npcs}

\begin{longtable}{|p{0.2\textwidth}|p{0.25\textwidth}|p{0.3\textwidth}|p{0.15\textwidth}|}
\hline
\textbf{Name} & \textbf{Role/Title} & \textbf{Motivation} & \textbf{Rumoured Quirk} \\
\hline
Sheikh Mansur & Oasis lord of the Three Wells & Keep his well's water rights against a rival clan & He has not tasted water in ten years – only coffee. \\
\hline
Fatima al-Wadi & Well-Judge of the central route & Settle disputes without bloodshed & She carries a dipper made from a date-palm leaf; it rots after each case. \\
\hline
Rashid & Caravan master of the Spice Road & Find a new route past a dried well & He has three wives, each from a different oasis. The wives do not speak. \\
\hline
Layla & Coffee-house poet & Sell verses for coin and secrets & Her poems are in code; the code is the price of coffee. \\
\hline
The Drowned Well-Keeper & A ghost of a well that went dry & Warn travellers of poisoned springs & She appears as a column of steam at dusk. She points at safe wells. \\
\hline
Harith al-Qalat & Last living survivor of Qalat al-Madina (immortal? cursed?) & Lift the curse on his city – or die trying & He drinks only rainwater. He has not aged in three centuries. He speaks Kuvani with a Fhara accent. \\
\hline
\textbf{Nerethi Whisperer} & \textbf{Psionic Agent of the Deep} & \textbf{Erase all debts and memories} & \textbf{Has no face; wears the skin of those they have erased.} \\
\hline
\end{longtable}

\subsection*{Names of the Fhara}
\label{sec:fhara-names}

Inspired by the desert peninsula: soft consonants, vowel-heavy, with clan prefixes (bin, bint, al-). Surnames often denote a well, a grove, or a trade.

\begin{longtable}{|p{0.25\textwidth}|p{0.25\textwidth}|p{0.2\textwidth}|p{0.2\textwidth}|}
\hline
\textbf{First Names (Male)} & \textbf{First Names (Female)} & \textbf{Surnames} & \textbf{Epithets} \\
\hline
Mansur & Fatima & al-Wadi & \textit{the Well-Born} \\
\hline
Rashid & Layla & bin Hassan & \textit{Coffee-Tongue} \\
\hline
Tariq & Amira & al-Sahel & \textit{the Ninth Cup} \\
\hline
Jamal & Nadia & Qasimi & \textit{Date-Father} \\
\hline
Hakim & Samira & al-Nur & \textit{the Dry-Handed} \\
\hline
Harith & Zubayda & al-Madina & \textit{the Drowned} \\
\hline
\end{longtable}

\paragraph*{(Well/Court/Bazaar)} A well-court with a stone bench and a judge's dipper; a coffee bazaar under canvas, roasting beans; a dried riverbed where old caravans are buried.

\section*{Spades — Places}
\label{sec:fhara-places}
\begin{enumerate}
\setcounter{enumi}{1}
\item \textbf{Well-Court} – A stone bench under a palm; a judge with a clay dipper.
  \hfill \textit{The dipper has a chip. The chip is from a sword that broke a well – and a trust.} \texttt{[JUSTICE]}
\item \textbf{Coffee Bazaar} – Stalls under canvas; the smell of roasting beans.
  \hfill \textit{The best coffee is served in a chipped cup. The chip is a signature.} \texttt{[TRADE]}
\item \textbf{Date Grove} – Rows of palms; the ground soft with fallen fruit.
  \hfill \textit{The fruit ferments at night. The drunk monkeys are the guards. They accept bribes of honey.} \texttt{[COMMUNITY]}
\item \textbf{Dry Riverbed (Old Caravan Route)} – A wide, sandy channel with bleached bones.
  \hfill \textit{The bones are not human. They are camels that died of thirst. Their ghosts still groan a warning.} \texttt{[TRAVEL]}
\item \textbf{Tent Court (Mobile)} – A carpet under a canvas roof; a judge with a scale.
  \hfill \textit{The scale weighs water, not silver. One side is always empty – for the next dispute.} \texttt{[JUSTICE]}
\item \textbf{Caravanserai of the Three Wells} – A walled courtyard with three different springs.
  \hfill \textit{The sweet well is for guests. The brackish well is for animals. The bitter well is for prisoners.} \texttt{[TRAVEL]}
\item \textbf{Salt Pan (Western Edge)} – A white flat where the caravans load salt.
  \hfill \textit{The salt is mixed with ancient tears. It tastes of old grief. Cooks use it sparingly.} \texttt{[TRADE]}
\item \textbf{Lighthouse of the Drowned} – A stone tower on the coast, never lit.
  \hfill \textit{The keeper is blind. He rings a bell when he hears ships. The ships are ghosts.} \texttt{[DEATH]}
\item \textbf{Oasis of the Ninth Palm} – A legend: a grove with nine palms, but the ninth is invisible.
  \hfill \textit{Drink from the ninth well, and you will forget your name for a day. Your friends will remind you.} \texttt{[LEGEND]}
\item[J] \textbf{Abandoned Well-Keeper's Hut} – A collapsed stone shelter with a dry basin.
  \hfill \textit{The basin is carved with a name. The name is the keeper's. He died of thirst – but the well is full.} \texttt{[DEATH]}
\item[Q] \textbf{Coffee Warehouse} – A cool, dark building with sacks of beans from a dozen oases.
  \hfill \textit{Each sack is stamped with a well's mark. The marks are family crests. The beans gossip.} \texttt{[TRADE]}
\item[K] \textbf{Well-Judge's Archive} – A cave filled with clay tablets. Each tablet is a water judgment.
  \hfill \textit{The tablets are baked hard. They are also porous. They weep when a judgment is reversed.} \texttt{[JUSTICE]}
\item[A] \textbf{The First Well} – A mythic spring at the centre of the peninsula, said to never dry.
  \hfill \textit{No one has found it. Those who claim they have are always lying – or dead. The dead are not lying.} \texttt{[LEGEND]}
\item[A+] \textbf{Qalat al-Madina (Sunken City)} – A field of sinkholes and half-submerged towers.
  \hfill \textit{The water is brackish. The ghosts are polite. They offer you a cup. The cup is empty. Drink it anyway.} \texttt{[RUIN]} \texttt{[DEATH]}
\end{enumerate}

\paragraph*{(Judge/Poet/Master)} Well-Judge with a clay dipper and a rotted palm-leaf; coffee-house poet with a stringed instrument; caravan master with a rope of camel hair.

\section*{Hearts — People \& Factions}
\label{sec:fhara-people}
\begin{enumerate}
\setcounter{enumi}{1}
\item \textbf{Well-Judge} – Grey-bearded, blind in one eye. He tastes the water before each ruling.
  \hfill \textit{``Salt means guilt,'' he says. ``Sweet means innocence. Bitter means appeal – and a cup of coffee.''} \texttt{[JUSTICE]}
\item \textbf{Coffee-House Poet} – A woman with hennaed hands and a silver bell on her ankle.
  \hfill \textit{She recites a verse. The verse is a price. The price is never in coin – it is a favor, a memory, a silence.} \texttt{[SOCIAL]}
\item \textbf{Caravan Master} – Leathery skin, a brass whistle on a thong.
  \hfill \textit{He can find water by the colour of the sand. He has never been wrong. He charges for the lesson.} \texttt{[TRAVEL]}
\item \textbf{Water-Clerk} – A young woman with a clay tablet and a reed pen.
  \hfill \textit{She records every bucket drawn. Her ink is water mixed with soot. A dry tablet means a dry well.} \texttt{[COMMUNITY]}
\item \textbf{Date Merchant} – Fat, cheerful, missing two fingers to a knife fight.
  \hfill \textit{He offers a date to every customer. The date is always from his own grove. He knows its story.} \texttt{[TRADE]}
\item \textbf{Oasis Lord (Sheikh)} – Robed in white, a bronze key on his belt.
  \hfill \textit{The key opens the sluice gate. He never lets it leave his sight. He sleeps with it under his tongue.} \texttt{[AUTHORITY]}
\item \textbf{Well-Spirit (Kulub)} – A shimmer of heat above the water at noon.
  \hfill \textit{It asks: ``Who drew without sharing?'' Answer truthfully, or the well will dry. The sharing is a story.} \texttt{[SPIRIT]}
\item \textbf{Camel Breeder} – A tall, silent man with a staff of acacia wood.
  \hfill \textit{He names each camel after a dead wife. The camels answer to the names. They never forget a slight.} \texttt{[COMMUNITY]}
\item \textbf{Salt Carrier} – A slave (illegal, but present in the deep desert). He wears a collar of braided rope.
  \hfill \textit{He is saving to buy his freedom. The price is a thousand measures of salt. His guild saves with him.} \texttt{[LABOR]}
\item[J] \textbf{The Drowned Well-Keeper's Ghost} – A woman in wet robes, standing at a dry well.
  \hfill \textit{She offers a cup of water. The cup is empty. Drink it anyway. The taste is memory.} \texttt{[DEATH]}
\item[Q] \textbf{Coffee Roaster} – A man with blackened hands and a brass drum.
  \hfill \textit{He roasts beans while reciting the genealogy of each grove. The beans remember. They also judge.} \texttt{[TRADE]}
\item[K] \textbf{The Ninth Palm's Guardian} – A hooded figure who appears at the legendary oasis.
  \hfill \textit{He asks: ``Why do you seek the invisible well?'' The answer determines if you see it. There is no wrong answer – only a price.} \texttt{[MYSTERY]}
\item[A] \textbf{The First Well-Keeper (Myth)} – A skeleton sitting on a stone bench, still holding a dipper.
  \hfill \textit{If you take the dipper, you become the new well-keeper. You cannot leave. The water is sweet. The silence is sweeter.} \texttt{[LEGEND]}
\item[A+] \textbf{Harith al-Qalat} – The immortal survivor of the Sunken City.
  \hfill \textit{He drinks only rainwater. He has not aged in three centuries. He knows where the Drowned Coin lies. He will not say.} \texttt{[LEGEND]} \texttt{[CURSED]}
\item[N] \textbf{Nerethi Mind-Thief} – A faceless figure wrapped in sand-coloured silk.
  \hfill \textit{They offer to take your pain. The price is the memory of who hurt you.} \texttt{[PSIONIC]}
\end{enumerate}

\paragraph*{(Dry/Coffee/Dispute)} A well is found dry – the keeper is dead; a coffee ritual is interrupted – the third cup is not poured; a water dispute erupts between two wells – the judge demands a bribe.

\section*{Clubs — Complications/Threats}
\label{sec:fhara-complications}
\begin{enumerate}
\setcounter{enumi}{1}
\item \textbf{Dry Well} – The oasis keeper is dead. Lose 1 Supply.
  \hfill \textit{The keeper's body is curled around the basin. He died of thirst. The well is full – but the spirit is angry.} \texttt{[DEATH]} \texttt{[SURVIVAL]}
\item \textbf{Coffee Ritual Interrupted} – A stranger refuses the third cup. The host draws a knife.
  \hfill \textit{The stranger is a rival clan member. The insult is a declaration of feud. The coffee is still hot.} \texttt{[SOCIAL]} \texttt{[COMBAT]}
\item \textbf{Water Dispute} – Two well-keepers claim the same spring. The judge demands a story.
  \hfill \textit{The story must be true. The judge will know if you lie. She has heard them all.} \texttt{[JUSTICE]} \texttt{[SOCIAL]}
\item \textbf{Date Blight} – A fungus attacks the groves. Prices triple. Starvation looms.
  \hfill \textit{The blight was spread by a merchant from the coast. The merchant's ship flew no flag. The dates weep.} \texttt{[DISASTER]} \texttt{[TRADE]}
\item \textbf{Sandstorm (Khamsin)} – Red dust obscures the path. Navigation DV +1, mark Fatigue.
  \hfill \textit{The dust carries whispers. The whispers are the names of dead caravaners. They are asking for directions.} \texttt{[WEATHER]}
\item \textbf{Customs Inspection (Forged Seal)} – A local governor demands papers. Yours are false.
  \hfill \textit{The seal is a copy. The forger is in the party. He is sweating. The governor is counting the drops.} \texttt{[AUTHORITY]} \texttt{[TRAP]}
\item \textbf{Elephant Stampede (Rare)} – A herd of stray elephants thunders through camp.
  \hfill \textit{The elephants are from a Dhaharan stable. They are escaped. They are looking for their mahout. He is hiding in your tent.} \texttt{[DISASTER]} \texttt{[COMBAT]}
\item \textbf{Mirage (False Well)} – A shimmer on the horizon. Lose 1 Supply chasing it.
  \hfill \textit{The mirage is a memory. The memory is of a well that dried a century ago. The ghost of the well drinks your water.} \texttt{[WEATHER]} \texttt{[TRAP]}
\item \textbf{Obligation Collector (Tulkani Memory-Cord)} – Someone recognises a knot you carry. They demand a story.
  \hfill \textit{The story must be about a broken promise. The knot tightens as you speak.} \texttt{[SOCIAL]} \texttt{[OBLIGATION]}
\item[J] \textbf{Well-Spirit Offended} – A traveller urinated in the basin. The well refuses water.
  \hfill \textit{The spirit demands a story. The story must be true. The well will remember.} \texttt{[SPIRIT]} \texttt{[SOCIAL]}
\item[Q] \textbf{Coffee House Fire} – A roasting drum explodes. The bazaar is burning.
  \hfill \textit{The fire was set by a rival oasis. The rival wants to monopolise the coffee trade. The smoke smells of cinnamon and spite.} \texttt{[DISASTER]} \texttt{[TRADE]}
\item[K] \textbf{Well-Judge's Verdict Forged} – A clay tablet is found altered. The judge is blamed.
  \hfill \textit{The forger is the judge's own clerk. The clerk owes a favour to the Salt Syndicate. The tablet weeps – but the tears are ink.} \texttt{[JUSTICE]} \texttt{[TRAP]}
\item[A] \textbf{The First Well Appears} – The mythic spring manifests. All who drink from it forget their broken promises.
  \hfill \textit{The forgetting is permanent. The sand shifts. The Hollow is pleased. The creditors are not.} \texttt{[LEGEND]} \texttt{[MAGIC]}
\item[A+] \textbf{Qalat al-Madina's Curse} – The Sunken City's ghosts rise and follow the party.
  \hfill \textit{They are polite. They offer water. The water is brackish. They will not leave until someone pays the Drowned Coin – or speaks the word the Omens whispered.} \texttt{[DEATH]} \texttt{[CURSE]}
\item[N] \textbf{Nerethi Memory Erasure} – A party member wakes up unable to recall a Bond or Obligation.
  \hfill \textit{The gap in their mind hums. The Nerethi are watching.} \texttt{[PSIONIC]} \texttt{[CURSE]}
\end{enumerate}

\paragraph*{(Water-share/Court/Dipper)} Water-share tablet (right to draw from a named well for a season); caravan court writ (neutral arbitration for one dispute); well-judge's dipper (one reroll on a negotiation involving water).

\section*{Diamonds — Rewards/Leverage}
\label{sec:fhara-rewards}
\begin{enumerate}
\setcounter{enumi}{1}
\item \textbf{Water-Share Tablet (Stewardship)} – A clay tile. Right to draw from a named well for one season.
  \hfill \textit{The tablet is a seal of trust. Break it, and the well's water becomes free for all – but the breaker becomes the new steward.} \texttt{[COMMUNITY]}
\item \textbf{Caravan Court Writ} – A palm-leaf document. Demands neutral arbitration for any dispute.
  \hfill \textit{The writ is sealed with coffee grounds. The seal is edible. Eat it to void the writ – and the dispute.} \texttt{[JUSTICE]}
\item \textbf{Well-Judge's Dipper} – A clay cup on a stick. Once, reroll a failed negotiation involving water.
  \hfill \textit{The dipper has a chip. The chip is from a sword that broke a well – and a trust. It knows a lie by taste.} \texttt{[SOCIAL]}
\item \textbf{Coffee Cup (Royal Seal)} – A chipped cup used by a Sheikh. Gain +1 Position in any Fhara court.
  \hfill \textit{The chip is a signature. No one will question it. The cup is still warm.} \texttt{[SOCIAL]}
\item \textbf{Date-Palm Seed} – A single seed from a sacred grove. Plant it to claim water rights.
  \hfill \textit{The seed grows into a palm within a year. The well-spirit will recognise you. The dates are a tithe.} \texttt{[COMMUNITY]}
\item \textbf{Oasis Lord's Key} – A bronze key to a sluice gate. Open any well's flow for one hour.
  \hfill \textit{The key is warm. It glows when a well is being poisoned. It burns the hand of a liar.} \texttt{[AUTHORITY]}
\item \textbf{Coffee Merchant's Record} – A record of every coffee sale for a decade. Contains names of spies.
  \hfill \textit{The record is written in disappearing ink. It reappears when heated. The spies are listed in the margins.} \texttt{[TRADE]} \texttt{[SECRET]}
\item \textbf{Salt Carrier's Rope Collar} – A braided rope that once held a slave. Present it to free a captive.
  \hfill \textit{The rope is old. It smells of sweat and freedom. One use. The freed slave will owe you a story.} \texttt{[COMMUNITY]}
\item \textbf{Well-Spirit's Boon} – A vial of water from a haunted well. Drink it to see the nearest spring.
  \hfill \textit{The water is cold. It tastes of copper. The vision lasts one hour. The spring may be guarded.} \texttt{[MAGIC]}
\item[J] \textbf{The Drowned Keeper's Cup} – An empty clay cup. Offer it to a dry well, and the dead keeper will fill it once.
  \hfill \textit{The cup weeps water. The water is salty. Drink it anyway. The dead keeper will ask for a story in return.} \texttt{[DEATH]}
\item[Q] \textbf{Coffee Roaster's Drum} – A brass drum. Roast beans in it, and the smoke forms words.
  \hfill \textit{The words are the future price of coffee. The price is always a warning. The smoke smells of regret.} \texttt{[MAGIC]} \texttt{[OMEN]}
\item[K] \textbf{The Ninth Palm's Date} – A single fruit from the invisible tree. Eat it to forget a single broken promise.
  \hfill \textit{The date is sweet. The forgetting is permanent. The well-spirit will still remember – but it will not collect.} \texttt{[MAGIC]}
\item[A] \textbf{The First Well's Water} – A flask from the mythic spring. Drink it to clear all obligations to Fhara patrons.
  \hfill \textit{The water tastes of nothing. The spirits will not forgive you for drinking it. But they will forget you owed.} \texttt{[LEGEND]} \texttt{[MAGIC]}
\item[A+] \textbf{Drowned Coin of Qalat al-Madina} – A silver piece minted in the lost capital.
  \hfill \textit{Present it to any Fhara well-spirit, and it will grant safe passage once. The coin was paid to a murderer. It still smells of blood – and of the Omens' whisper.} \texttt{[LEGEND]} \texttt{[OBLIGATION]}
\item[N] \textbf{Nerethi Silence} – A psionic token that blocks all Obligation tracking for one session.
  \hfill \textit{You feel light. You feel empty. You do not remember why you came here.} \texttt{[PSIONIC]} \texttt{[LEVERAGE]}
\end{enumerate}

\section*{Quick use notes}
\label{sec:fhara-quick-use}
\begin{itemize}
\item Draw until you have all four suits: Spade = place, Heart = actor, Club = pressure, Diamond = leverage. Highest rank sets the main Timer (2--5 \(\rightarrow\) 4, 6--10 \(\rightarrow\) 6, J/Q/K \(\rightarrow\) 8, A \(\rightarrow\) 10).
\item Diamonds are codified outcomes (tablets, writs, dippers) that change position rather than call for a roll.
\item If any A appears, echo \textbf{water-omens} – a dry well that weeps, coffee that brews itself, a date that bleeds.
\item If the party visits Qalat al-Madina, treat it as a Class IV Anchor site (see Saikou Ira's compendium). The ghosts are not hostile, but they are very, very patient – and they remember the Omens' warning.
\end{itemize}

\subsection*{Additional Features (Cultural Mechanics)}
\label{sec:fhara-mechanics}

\begin{itemize}
\item \textbf{Three Cups, One Contract:} Fhara negotiations follow a rigid ritual. First cup: welcome and genealogy. Second cup: bargaining and threats disguised as compliments. Third cup: sealing the deal. Refusing the third cup voids the negotiation and insults the host. Accepting the third cup without intending to honour the deal is a capital crime – the liar is tied to a date palm until the sun sets.
\item \textbf{Water-Sharing Tablets:} Water rights are recorded on clay tablets, stored in the Well-Judge's archive. A tablet can be split, traded, or shared. The tablet is the well's memory; break it, and the water becomes unclaimed – but the breaker becomes the new steward.
\item \textbf{Well-Spirits (Kulub):} Each well has a spirit, usually benign but watchful. Offerings of dates, coffee, or a true story keep it calm. A well that has witnessed a murder may turn bitter; a well that has been poisoned may turn sentient. A sentient well is a disaster – it follows those who wronged it.
\item \textbf{The Drowned Coin and the Omen-Whisper:} Silver pieces from Qalat al-Madina still circulate. A Drowned Coin can be used to invoke the Water-Right of the Lost Capital – one free, unquestioned draw from any well in Fhara territory. The well-spirit will remember, and it will ask for a story in return. But if you listen carefully while holding the coin, you may hear the Omens' whisper: \textit{``Remember.''}
\item \textbf{Local Patrons (Syncretic Names):}
  \begin{itemize}
    \item \textbf{The Thirst} – Khemesh. \textit{The hunger beneath the sand; the pressure that drinks wells dry.}
    \item \textbf{The Cup Bearer} – The Witness. \textit{The third cup that never lies; the contract sealed in coffee grounds.}
    \item \textbf{The Oasis Father} – The Traveler. \textit{The road that leads to water; the guide who never drinks.}
    \item \textbf{The Date Mother} – Inaea. \textit{The web of groves; the shared harvest that binds families.}
    \item \textbf{The Sunken King} – Varnek Karn. \textit{The ghost of Qalat al-Madina; the negotiator of secrets that outlast cities.}
  \end{itemize}
\end{itemize}

\subsection*{Lore Echoes (Cross-Expansion Hooks)}
\label{sec:fhara-lore}

\begin{itemize}
  \item \textbf{The Omen-Whisper (Inaea, Ikasha, Isoka):} In the deepest wells of the Sunken City, a traveller may hear a voice speaking a single word: \textit{``Remember.''} The voice is not a ghost – it is a fragment of the three Omens' warning, left behind when they commanded the Ykrul and Kuvani to bury the ancient knowledge. A character who drinks from the Cursed Spring and survives may hear the full whisper – a single sentence in a language older than the empire. The GM may reveal a clue about the forbidden vault.
  \item \textbf{Coffee and the Covenant (Mykkiel):} Fhara coffee houses are neutral ground under Covenant law. A well-judge's verdict cannot be appealed outside the oasis. Some Covenant judges are trained in the three-cup ritual – they can tell a liar by the way they hold the cup.
  \item \textbf{Water Tablets and the Aeler Ledger (Ngomebe):} Aeler engineers have been hired to map the underground aquifers. Their breath-books record water pressure and flow. The Fhara see this as an intrusion – or an opportunity for better wells. A shared aquifer is a shared treaty.
  \item \textbf{Date Groves and the Ukwe (Sekogo):} The Ukwe's roots do not reach the dry peninsula, but the Fhara have their own Speaking Palms – trees that remember every date harvested. A date from a sacred grove can witness an oath as surely as a Lethai-ar cord. The tree will not lie.
  \item \textbf{The First Well and the Hollow (Eastern Realms):} The mythic spring that never dries is said to be the Hollow's original well – before it learned to thirst. Aelaerem hedge-keepers sometimes leave offerings at dried wells, hoping to appease the Hollow with water. The Hollow does not drink. It counts.
  \item \textbf{The Sunken City and the Kuvani (Eastern Ykrul):} The Kuvani and Eastern Ykrul never forgot Qalat al-Madina. Some say they still patrol its ruins to prevent the Fhara from reclaiming it – and to ensure the Omens' seal remains unbroken. Others say the curse was their doing, and they hold the only key to lifting it: a bronze paiza stamped with the city's original water treaty, given to them by Inaea herself.
  \item \textbf{The Nerethi and the Void (Psionics):} The Nerethi claim the Hollow is a lie told by creditors. They seek to erase the Ninth Taboo entirely. If they succeed, the Sunken City will rise, and the debt of the world will be wiped clean – along with its history.
\end{itemize}

\subsection*{Fhara Superstitions (burn for +1 die once)}
\begin{itemize}
  \item \textbf{Bread to the Basin:} Crumble a crust into the well. Ignore the first \emph{Dry Well} penalty. \texttt{[SUPERSTITION]}
  \item \textbf{Knot the Dipper:} Tie a single date-palm fibre around the dipper's handle. Cancel \emph{Coffee Ritual Interrupted} or take --1 die on your next negotiation (your choice). \texttt{[SUPERSTITION]}
  \item \textbf{Name to the Well:} Whisper a dead well-keeper's name into the water. Gain +1 Effect when haggling over water rights this scene, but mark Obligation to the Well-Spirit. \texttt{[SUPERSTITION]}
  \item \textbf{Coin to the Ghost:} Toss a Drowned Coin into a dry well. The ghost of Qalat al-Madina will guide you to water once – and whisper the first syllable of the Omens' warning. \texttt{[SUPERSTITION]}
\end{itemize}

\subsection*{Thematic SB Spend Table}
\label{sec:fhara-sb}

\paragraph{Minor Complications (1 SB)}
\begin{itemize}
\item \textbf{Exposure:} Your actions draw unwanted attention from \textbf{Well-Judges or Water-Clerks}.
\item \textbf{Noise:} Sounds of your actions alert nearby \textbf{coffee-house guards or date merchants}.
\item \textbf{Trace:} Evidence of your passage marks your route for \textbf{caravan masters or well-spirits}.
\item \textbf{Delay:} A brief but meaningful setback costs you \textbf{time or favourable water access}.
\item \textbf{Supply Strain:} Mark +1 segment on a relevant \textbf{resource timer}.
\end{itemize}

\paragraph{Moderate Setbacks (2 SB)}
\begin{itemize}
\item \textbf{Alarm Raised:} \textbf{Local Well-Judge or Oasis Lord} becomes aware and begins responding.
\item \textbf{Position Lost:} You lose advantageous ground/cover/stealth due to \textbf{sandstorm or mirage}.
\item \textbf{Foe Appears:} A \textbf{rival oasis guard or obligation collector} arrives on scene.
\item \textbf{Gear Trouble:} A piece of equipment becomes \textbf{Compromised/Neglected}.
\item \textbf{Lock/Barrier:} A simple obstacle now requires a test to overcome.
\end{itemize}

\paragraph{Serious Trouble (3 SB)}
\begin{itemize}
\item \textbf{Reinforcements:} Additional \textbf{well-guards, coffee-house enforcers, or camel riders} arrive.
\item \textbf{Key Gear Breaks:} A crucial tool/weapon becomes temporarily unusable.
\item \textbf{Major Twist:} The situation fundamentally changes – \textbf{well dries / coffee warehouse burns / water tablet breaks}.
\item \textbf{Rail Tick:} Advance a relevant campaign/front timer by 1 segment.
\item \textbf{Condition Applied:} Mark \textbf{Fatigue 1/Harm 1/Condition} appropriate to fiction.
\end{itemize}

\paragraph{Major Turns (4+ SB)}
\begin{itemize}
\item \textbf{Trap Springs:} A prepared danger activates with full effect.
\item \textbf{Authority Arrival:} \textbf{High Well-Judge, Oasis Sheikh, or First Well-Keeper's ghost} intervenes.
\item \textbf{Scene Shift:} The environment changes dramatically – \textbf{well collapses / oasis floods / date grove burns}.
\item \textbf{Patron Omen:} Divine/arcane forces take notice – \textbf{omen appears / blessing lost / curse manifests}.
\item \textbf{Narrative Pivot:} The story takes an unexpected turn that reframes objectives.
\end{itemize}

\paragraph{Region-Specific SB Options}
\begin{itemize}
\item \textbf{Fhara (Water Law):} Well-spirits ripple the surface, clay tablets glow, a dipper's chip grows larger.
\item \textbf{Fhara (Coffee Ritual):} Beans roast themselves, cups overflow without water, the third cup pours black smoke.
\item \textbf{Fhara (Desert):} Sand shifts to reveal old bones, mirages speak, a dry well whispers a name.
\item \textbf{Fhara (Sunken City):} A drowned tablet floats up from a sinkhole, still legible. Its judgment is from a dead judge. It is still binding – and it mentions the Omens.
\item \textbf{Fhara (Nerethi):} The air hums. A party member forgets why they drew their weapon. A Bond vanishes from the sheet.
\end{itemize}

\subsection*{Pursuit Ladder (Near / Mid / Far)}
\begin{itemize}
\item \textbf{Advance/Withdraw (Wits + Survival):} On a hit, shift one band. On a strong hit, also impose \emph{Sand Blind}.
\item \textbf{Cut the Bargain (Presence + Sway):} On a hit, freeze range for one exchange; if a water-share tablet is in your possession, +1 die.
\item \textbf{Weather the Khamsin (Resolve + Endurance):} On a hit in \emph{Sandstorm}, you pick who stays on their feet; on a miss, the GM does.
\end{itemize}

\subsection*{Boss Seeds (Water \& Sand)}
\begin{itemize}
\item \textbf{The Drought-Maker (Nine Dry Wells)} – \emph{Phases}: Well Poisonings \(\rightarrow\) Orchestrated Famine \(\rightarrow\) Oasis Coup. \emph{Cracks}: find a hidden spring, expose a false tablet, survive a season without water.
\item \textbf{The Coffee-Price (Third Cup Cursed)} – \emph{Phases}: Ritual Interruptions \(\rightarrow\) Procession of Empty Cups \(\rightarrow\) Bargain Baptism of a Merchant. \emph{Cracks}: relic coffee bean, survivor testimony, the third cup turning sweet again.
\item \textbf{The Sunken King (Ghost of Qalat al-Madina)} – \emph{Phases}: Drowned Archive \(\rightarrow\) Curse of the Ninth Well \(\rightarrow\) The Coin's Reckoning. \emph{Cracks}: return the Drowned Coin to his skeletal hand, recite the lost water-treaty from a preserved tablet, or convince a living Fhara sheikh to publicly forgive the Kuvani – and to speak the Omen-Whisper aloud.
\item \textbf{The Nerethi Hive-Mind} – \emph{Phases}: Memory Theft \(\rightarrow\) Psionic Intrusion \(\rightarrow\) Total Erasure. \emph{Cracks}: protect a Bond, recover a lost memory, force the Hive to owe a debt.
\end{itemize}

\begin{tcolorbox}[colback=black!3,colframe=black!40!white,title={Patronage \& Power}]
Among the Fhara, power flows through water, coffee, and the weight of a shared bargain. The Well-Judge's word is law, the Oasis Lord controls the sluice, and the Caravan Master commands the desert. But no one rules alone – every decision is a negotiation, every negotiation a ritual, every ritual a shared trust.

\textbf{For the GM:}  
Power in Fhara society is measured in well-depth, coffee cups, and clay tablets. Patronage often takes the form of water-share rights, court writs, or a judge's dipper – each tied to a specific oasis. To emphasise community and reciprocity rather than obligation:
\begin{itemize}
\item Tie rewards to physical tokens (tablets, dippers, keys) that represent shared access rather than obligation. Breaking a token is breaking a trust, not a contract.
\item Let rival oases issue conflicting water claims – resolving them requires a community feast or a joint coffee ritual, not a payment.
\item Use the well-courts, coffee bazaars, and date groves as arenas for social contests, where a cup of coffee is a shared moment and a dry well is a shared tragedy.
\item Introduce the Omen-Whisper as a recurring mystery. The party may find fragments of the Omens' warning carved into ancient well-stones – each fragment a clue to the forbidden knowledge buried beneath Qalat al-Madina.
\item \textbf{Nerethi Contrast:} Use the Nerethi to challenge the party's attachment to their Bonds and Obligations. Is a debt worth remembering if it causes pain?
\end{itemize}
In Fhara society, the sand forgets nothing – and the well remembers every cup shared.
\end{tcolorbox}

% Index entries
\index{Fhara|see{Oasis Courts}}
\index{Theme generator}
\index{Oasis Courts generator}
\index{Places!Fhara}
\index{People!Fhara}
\index{Complications!Fhara}
\index{Rewards!Fhara}
\index{Well-Judge}
\index{Water-Clerk}
\index{Qalat al-Madina}
\index{Sunken City}
\index{Omen-Whisper}
\index{Inaea}
\index{Ikasha}
\index{Isoka}
\index{Eastern Ykrul}
\index{Kuvani}
\index{Local Patrons!Fhara}
\index{Nerethi}
\subsection*{The Narethi — The Unburied of the Red Waste}
\label{subsec:narethi}
\index{Narethi}\index{Fhara!Narethi}

\emph{``The sun remembers what the sand forgets. We remember both. That is our curse. That is our strength.''}

Deep in the Red Waste, where the Fharan oases give way to cracked salt flats and the bones of failed cities, dwell the Narethi. They are not Fhara — they were never Fhara — but they have lived on the edge of Fhara's desert for so long that the water-clerks have learned to leave offerings at the boundary stones and ask no questions.

The Narethi are felid — cat-people, in the crude speech of lowland merchants — with tufted ears, slit-pupiled eyes that narrow to threads in the sun, and a proud, slow way of moving that suggests they are never quite where you expect them to be. Their fur ranges from the pale cream of dried salt to the deep rust of the Red Waste's iron-rich sands. Their voices carry a harmonic resonance, two pitches at once, as if they are always speaking in duet with a shadow of themselves.

What you notice first: The way they never blink in unison. The scent of myrrh and old copper. The silence when they stop speaking — it is not empty; something is listening.
The one rule that kills newcomers: Never ask a Narethi about the ninth syllable. They will not answer. They will not forget that you asked.

\paragraph{Where They Live}
The Narethi have no cities, no wells, no date groves. They dwell in the \textbf{Sunken Courts} — a network of subterranean chambers carved from ancient salt and basalt, hidden beneath the cracked flats of the eastern Red Waste. The entrance is a sinkhole that smells of ozone and old copper. Visitors who are not Narethi are blindfolded and led through nine switchbacks before the tunnels open into the first Court.

The Courts are warm, lit by phosphorescent crystals that the Narethi claim are the teeth of a god who died before the first Fhara well was dug. The air is still and dry, scented with myrrh, old blood, and the faint, sweet smell of mummification spices.

\paragraph{Origins (What the Narethi Believe)}
The Narethi do not speak of their origins to outsiders. When asked, they say: \textit{``We were the ones who did not burn. The world turned to glass. We did not turn with it. We waited. We are still waiting.''}

Their folklore tells of a \textbf{First Scouring} — a war in the sky before the first human king drew breath. They do not name the combatants. They say only that \textit{three stars argued, one dreafellmed, and one of the fallen's children tried to swallow the sun}. The Narethi chose the wrong side — or the right side, depending on which elder you ask — and were exiled to the wastes as punishment. Their task: to guard a \textbf{sleeping thing} buried beneath the salt, to sing it lullabies, and to ensure it never wakes.

They do not remember the war clearly. Their memories are echoes, passed down in harmonic chants that lose one syllable each generation. The oldest chant is now only eight syllables long. The ninth syllable was the name of the thing they guard.

\paragraph{The Unburied — The Mummified One}
At the heart of the Sunken Courts, in a chamber cooled by no draft and lit by no crystal, lies the \textbf{Unburied}. It is a Narethi — or was, once — wrapped in linen and salt, its fur preserved, its eyes sewn shut with silver wire, its mouth open in a silent scream that has not stopped for millennia.

The Narethi do not worship the Unburied. They attend to it. They say it was the first of them, the one who refused to burn, and that its mind still dreams beneath the salt. Its dreams are psionic — they leak into the Courts, shaping the crystals, influencing the Narethi's thoughts, and granting them their strange mental gifts.

The Unburied does not speak. It \textit{resonates}. Narethi who meditate in its chamber report hearing a single word, repeated endlessly, in a voice that sounds like their own: \textit{``Remember. Remember. Remember.''}

\paragraph{The Narethi's Psionic Gifts}
All Narethi are psions, though they do not use that word. They call their abilities \textit{the Resonance} — an echo of the Unburied's eternal dream. A Narethi's psionics are instinctive, shaped by meditation and the harmonic chants they learn as children.

\textbf{Narethi Psionic Traits (Optional Player Character)}
\begin{itemize}
    \item \textbf{Natural Telepathy:} You may communicate wordlessly with any willing creature within Near range. This is not language — it is emotion, image, and intent.
    \item \textbf{Resonance Sense:} Once per session, you may touch a surface (stone, salt, metal) and ask the GM one yes/no question about an event that occurred within Near range in the past day. The resonance lingers.
    \item \textbf{Disciplined Mind:} You gain +1 die to resist psychic assault or mental intrusion. Your thoughts are your own — and the Unburied's.
    \item \textbf{Night-Eyes:} Narethi eyes reflect light like a cat's. In darkness, you see as if in dim light. This is natural, not magical.
\end{itemize}

\paragraph{The Resonance Price}
The Narethi's psionic gifts are not free. Each time a Narethi uses a significant psionic ability (a Telepathy deep read, a Clairvoyance scry, a Psychic Assault), they mark a segment on their \textbf{Resonance Leash [4]}. When the leash fills, the Unburied's dream bleeds through — the Narethi experiences a waking vision of the First Scouring. The GM describes a fragment: a city of black glass, a sky on fire, a scream that has no end. The Narethi marks 1 Fatigue and the leash resets.

\paragraph{Society and Structure}
The Narethi are ruled by a \textbf{Father of Resonance} — an elder who has spent the most time in meditation with the Unburied. She does not command; she interprets the dream. Her word is law because the dream has foreseen the consequence of disobedience. Below her are \textbf{Chanters} (priests and lore-keepers), \textbf{Watchers} (scouts and warriors who patrol the boundaries of the Red Waste), and \textbf{Silent Ones} (Narethi who have taken a vow of non-speech to better hear the Resonance).

Their society is rigid, hierarchical, and secretive. Outsiders are not trusted. A Narethi who leaves the Courts is considered \textit{unbound} — they may return, but they will be questioned about everything they saw and heard. Some unbound Narethi never return. The Father says the dream does not speak of them.

\textbf{Factions Within the Courts:}
\begin{itemize}
    \item \textbf{The Listeners} — Believe the Unburied is trying to wake, and that the Narethi should help it. They study ancient chants, seeking the lost ninth syllable.
    \item \textbf{The Mutes} — Believe the Unburied should never wake. They have taken vows of absolute silence, hoping that if no one speaks, the dream will fade.
    \item \textbf{The Unbound} — Narethi who have left the Courts and lived among surface-dwellers. They are viewed with suspicion, but they are also the Narethi's only source of news about the outside world.
\end{itemize}

\paragraph{Relationship with the Fhara}
The Fhara know of the Narethi, but they do not speak of them. At the boundary stones that mark the edge of the Red Waste, well-judges leave offerings of salt and dried dates twice a year. The Narethi take them. No one meets.

When a Fhara caravan strays too close to the Sunken Courts, it may return with no memory of the detour, its water-skins refilled, its camels calm. The Narethi do not harm — they erase. A traveler who has been touched by the Resonance may wake with a single new knot in their memory-cord, tied in a pattern no Tulkani recognizes.

The Fhara call the Narethi \textit{the Unremembered}. It is not an insult. It is a warning.

\paragraph{The Secret the Mother Keeps}
The Narethi do not know why they guard the Unburied. They only know that if it wakes, the Red Waste will spread. The sand will turn to glass. The sun will not set for nine days.

What the Father of Resonance knows — and does not tell even the Chanters — is that the Unburied is not a prisoner. It is a \textit{lock}. Beneath it, buried deeper than the Courts, deeper than the salt, deeper than the bones of the First Scouring, something else sleeps. The Unburied's dream is the only thing keeping it asleep. If the Unburied ever stops dreaming, the thing beneath will remember what it was.

The Father does not know what that thing is. She only knows that the dream's single word — \textit{``Remember''} — is not a command to the Narethi. It is a command to the thing below. \textit{Remember what you are. Remember what you did. Remember why you must never wake.}

\clearpage
\subsubsection*{Narethi d6 Mini-Generators}

\paragraph{Narethi Places (Spades) — d6}
\begin{enumerate}
    \item \textbf{The Sunken Courts (Seat of Resonance)} — A subterranean complex of salt-crusted chambers lit by phosphorescent crystals. The crystals hum when the Unburied dreams deeply. Visitors are blindfolded through nine switchbacks. [PSIONIC] [SECRET]
    \item \textbf{The Chamber of the Unburied (Heart)} — A sealed room cooled by no draft, lit by no crystal. The mummified Narethi lies on a bed of salt and linen, eyes sewn shut with silver wire. The air tastes of copper and old grief. [DEATH] [PSIONIC]
    \item \textbf{The Hall of Lost Syllables (Archive)} — Walls carved with harmonic chants. The oldest chant has eight syllables; the ninth space is empty, polished smooth by generations of fingers. The stone is warm. [MEMORY] [RIDDLE]
    \item \textbf{The Resonant Well (Psionic Pool)} — A basin of still, dark water. Touch it, and you hear echoes of the Unburied's dream. Narethi meditate here to attune their Resonance. The water never ripples. [PSIONIC] [OMEN]
    \item \textbf{The Watchers' Perch (Boundary)} — A salt-crusted rock outcropping overlooking the Red Waste. Watchers sit here in silence, scanning for intruders. The perch is marked with nine carved eye symbols — the ninth is half-erased. [WATCH] [STEALTH]
    \item \textbf{The Salt Pan of Forgotten Names (Exile Ground)} — A cracked white flat where unbound Narethi who cannot return are said to walk. Their footprints appear at dawn and vanish by noon. Do not follow them. The salt tastes of tears. [DEATH] [MYSTERY]
\end{enumerate}

\paragraph{Narethi People \& Factions (Hearts) — d6}
\begin{enumerate}
    \item \textbf{Father of Resonance (Ruler)} — An elder with fur gone white, eyes that reflect no light. She speaks only to interpret the dream. Her voice is a whisper that carries through the Courts. She has not left the Chamber in twenty years. [AUTHORITY] [SECRET]
    \item \textbf{Listener Chanter (Priest)} — A Narethi with a shaved head, harmonic voice, and a cord of nine knots. She seeks the lost ninth syllable. She will trade a resonance secret for a story about a dream you cannot forget. [PSIONIC] [LORE]
    \item \textbf{Mute Sentinel (Silent One)} — A warrior who has taken a vow of silence. She communicates with hand signs and a small clay slate. Her silence is so complete that footsteps make no sound. She has not spoken in nine years. [STEALTH] [COMBAT]
    \item \textbf{Unbound Exile (Outsider)} — A Narethi who left the Courts and now works as a caravan scout. She wears a hood to hide her eyes. She will guide travelers past the Red Waste — but she asks in return for news of the Courts. She cannot return. [TRAVEL] [TRAGEDY]
    \item \textbf{The Dream-Echo (Haunted)} — A Narethi who spent too long in the Resonant Well. The Unburied's dream now bleeds into his waking hours. He speaks in two voices, one his own, one a whisper of ancient glass. He is looking for someone to cut the resonance. [PSIONIC] [CURSE]
    \item \textbf{The Keeper of Teeth (Crystal-Minder)} — An old, blind Narethi who tends the phosphorescent crystals. She claims the crystals are growing — that they were not here when she was young. She is afraid of what they will become. [CRAFT] [OMEN]
\end{enumerate}

\paragraph{Narethi Complications (Clubs) — d6}
\begin{enumerate}
    \item \textbf{Resonance Bleed} — The Unburied's dream intensifies. All Narethi within the Courts mark 1 Fatigue. Visitors hear whispers in a language they almost understand. The crystals glow brighter. [PSIONIC] [FEAR]
    \item \textbf{The Ninth Syllable Spoken} — Someone speaks the lost syllable aloud — by accident, or by design. The Unburied's mouth twitches. The ground trembles. Start a Waking Dream [6] timer. [DEATH] [APOCALYPSE]
    \item \textbf{Unbound Hunted} — A Narethi exile is being pursued by Watchers. The exile begs the party for passage across the border. Helping them earns a Narethi's gratitude — and the Mother's attention. [STEALTH] [POLITICS]
    \item \textbf{The Crystals Sing} — The phosphorescent crystals emit a harmonic tone that induces waking visions. Anyone who hears it must test Spirit+Resolve (DV 4) or experience a fragment of the First Scouring. Mark 1 Fatigue on failure. [PSIONIC] [OMEN]
    \item \textbf{Mute's Vow Broken} — A Silent One speaks for the first time in years. She speaks a warning: something is waking beneath the salt. The other Narethi do not believe her. She will not speak again. [TRUTH] [TRAGEDY]
    \item \textbf{The Salt Pan Walks} — Footprints in the Salt Pan of Forgotten Names begin moving toward the Courts. The dead unbound are coming home. They are not hostile — they are curious. That is worse. [DEATH] [FEAR]
\end{enumerate}

\paragraph{Narethi Rewards \& Leverage (Diamonds) — d6}
\begin{enumerate}
    \item \textbf{Resonance Shard (Crystal Fragment)} — A piece of phosphorescent crystal from the Courts. Once per session, touch it to gain +1 die on a psionic roll. The shard then dims for a scene. The Unburied knows you used it. [PSIONIC] [LEVERAGE]
    \item \textbf{Mute's Clay Slate (Witness Token)} — A small clay tablet carried by a Silent One. Present it to any Narethi to demand a single truthful answer. The slate then cracks. The Mute who gave it will not speak of it. [SOCIAL] [TRUTH]
    \item \textbf{Listener's Cord (Nine Knots)} — A knotted cord with eight tight knots and one loose. Untie a knot to ask the GM one yes/no question about the Unburied's dream. The knot reties itself after use — but the loose knot tightens slightly. [MEMORY] [RIDDLE]
    \item \textbf{Watcher's Salt Pouch (Boundary Offering)} — A leather pouch of salt blessed by the Watchers. Scatter it to create a temporary threshold that spirits and psionic projections cannot cross. Lasts one scene. The salt tastes of iron. [WARD] [PROTECTION]
    \item \textbf{Unbound's Map (Hidden Route)} — A scrap of hide marked with charcoal lines. Shows a safe route through the Red Waste that avoids both Fhara patrols and Narethi Watchers. The map was drawn by an exile who never returned. [TRAVEL] [STEALTH]
    \item \textbf{The Dream-Echo's Confession (Whisper)} — A single word whispered by a Dream-Echo before he went silent. The word is a fragment of the ninth syllable. Speaking it aloud to a Narethi grants +1 Position in any negotiation with them — but the Unburied will remember your voice. [LEGEND] [DANGER]
\end{enumerate}

\paragraph{Adventure Hooks (d6 Quick Starts)}
\begin{enumerate}
    \item \textbf{The Lost Unbound.} A Narethi who left the Courts a decade ago has resurfaced — as a mercenary captain in Ecktoria. The Father of Resonance sends the party to retrieve them, or to learn why they refuse to return. The Unbound has a fragment of the ninth syllable tattooed on their fur.
    \item \textbf{The Resonance Bleed.} A Fhara well has begun to taste of copper and ozone. Those who drink from it dream of black glass and a screaming sky. The Narethi say the Unburied is restless. Something has disturbed its slumber — and the thing beneath is stirring.
    \item \textbf{The Silent One's Plea.} A mute Narethi appears at a caravanserai, carrying a clay tablet carved with a single symbol. The symbol is not Narethi — it is older. The party must deliver the tablet to the Archivolt in Thepyrgos, where a Lethai-thora keeper recognizes it as a fragment of the First Scouring. The tablet begins to warm.
    \item \textbf{The Ninth Syllable.} A Tulkani story-singer has begun to hum a chant that is not Tulkani — a chant with nine syllables, not eight. The Narethi Listeners want to hire the party to capture the singer. The Mutes want the singer silenced — permanently. The singer does not know where the chant came from.
    \item \textbf{The Salt Pan Pilgrims.} A group of unbound Narethi exiles is walking toward the Courts from the Salt Pan. They have been dead for years. They are not hostile — they are singing a lullaby. The Father wants to know why they have returned. The party must follow them without being seen.
    \item \textbf{The Crystal's Hunger.} The phosphorescent crystals in the Courts have begun to grow toward the surface. A Watcher reports that they pulse in time with the Unburied's dream. The Keeper of Teeth says they are feeding on something — and that something is almost gone.
\end{itemize}

\paragraph{Player Takeaway}
The Narethi are not a character class or a mechanical package. They are a story — a people shaped by a burden they no longer remember, guarding a secret they cannot name. A Narethi player character is \textit{unbound}: they have left the Courts, for reasons of their own. They may be a spy, an exile, a pilgrim, or simply someone who could not bear the silence.

Work with your GM to determine:
\begin{itemize}
    \item Why did you leave the Sunken Courts?
    \item Do you still feel the Resonance? (If yes, you track a Resonance Leash. If no, you have lost some of your psionic gifts — but you are truly free.)
    \item What do you remember of the Unburied's dream?
    \item Who in the Courts wants you to return — and who wants you dead?
\end{itemize}

The Narethi are few. The Courts are quiet. The dream continues.

\noindent \textit{``The ninth syllable is not a sound. It is the space between heartbeats. The Unburied has been counting for a very long time.''}